\subsection{Tuple Syntax and Structural Role}

Tuples share list sequence operations for reading but omit \texttt{append}, \texttt{sort}, and other resize/mutate APIs. \texttt{count} and \texttt{index} remain because they inspect without altering slot bindings.

\subsubsection{Tuple Unpacking: Decomposing Fixed-Length Structures into Multiple Names}

Unpacking assignment \texttt{a, b, c = iterable} is lowered by the compiler to stack operations: the right-hand iterable is evaluated, then \texttt{UNPACK\_SEQUENCE} (with arity $n$) pops $n$ values and binds local names. Arity mismatch (\texttt{ValueError}: too many / too few values) is detected when the actual length differs from the target count---effectively a stack-shape trap at runtime. Extended unpacking with \texttt{*} routes surplus elements into a \texttt{list}, not a \texttt{tuple}.

The swap idiom \texttt{a, b = b, a} relies on tuple packing on the right (anonymous tuple construction on the stack) before left-hand unpacking---no temporary name required. Iteration \texttt{for x, y in pairs} applies the same unpack machinery each loop iteration.

\subsubsection{Single-Element Tuple Syntax: Why \texttt{(x,)} Is a Tuple but \texttt{(x)} Is Not}

Parentheses group expressions; commas construct tuples. \texttt{(42)} is an \texttt{int}; \texttt{(42,)} is a one-element \texttt{tuple}. The trailing comma is the tuple marker. Confusing a one-character \texttt{str} with a one-element tuple of that string changes \texttt{len} and structural semantics for APIs expecting tuple arity.
